\subsection{A Framework for Direct and Indirect Genetic Support}

\begin{figure}[t!]
    \centering
    {{canva:3}}
    \caption{Conceptual overview of the PIGEAN framework. (A) Disease-associated genes identified through GWAS or Mendelian databases. (B) Disease-relevant genes share functional annotations with associated genes. (C) Graphical model of PIGEAN. (D) Genes positioned along direct and indirect support axes. (E) Downstream applications including drug target prioritization and gene-trait relevance mapping.}
    \label{fig:pigean-method}
\end{figure}

We define a \textit{disease-associated gene} as one carrying genetic variants that reach statistical significance in GWAS or that has been curated in a Mendelian database. We define a \textit{disease-relevant gene} as one for which perturbing its activity can causally influence disease risk or progression. Disease-associated genes are a subset of disease-relevant genes, and our goal is to quantify evidence for both.

To estimate direct genetic support, we identified independent association signals via clumping and assigned signals to genes using a probabilistic nearest-gene approach that accounts for gene density and distance \cite{dornbos_evaluating_2022} (Methods). To estimate indirect genetic support, we developed PIGEAN (Priors Inferred from GEne ANnotations), a Bayesian framework in which functional annotations modulate the prior probability of gene relevance, with annotation effect sizes learned jointly from genetic data (Methods). Genes sharing annotations with associated genes accumulate indirect support even when they lack direct associations. Both direct and indirect support are expressed as Bayes factors on a common scale, and indirect support exceeding 1 indicates evidence for disease relevance beyond baseline.

A key feature of PIGEAN is a spike-and-slab prior on geneset effect sizes that strictly prunes redundant or trait-irrelevant annotations. This allows us to include large heterogeneous annotation libraries representing broad swaths of current biological knowledge without inflating false positives. Genetic associations serve as the anchor: only annotations enriched among associated genes receive non-zero effects, while irrelevant or redundant annotations are regularized to zero. This regularization means that indirect support for a gene reflects only the specific annotations that are demonstrably enriched among associated genes for the trait in question, rather than a nonspecific accumulation of biological knowledge. As a consequence, indirect support provides a principled probabilistic justification for several practices the field currently performs heuristically, including prioritizing genes at GWAS loci based on prior biological knowledge, interpreting suggestive signals below genome-wide significance, and calibrating the relative informativeness of different non-genetic evidence sources (Supplementary Information).

We applied PIGEAN across {{NUM_COMPLEX_TRAITS:,}} complex traits and {{NUM_RARE_DISEASE_GROUPS:,}} rare disease groups using {{NUM_TOTAL_GENE_ANNOTATIONS:,}} gene annotations spanning knockout mouse phenotypes ({{NUM_MPO_TERMS:,}} MPO terms), curated pathways ({{NUM_MSIGDB_GENESETS:,}} MSigDB gene sets), literature co-occurrence ({{NUM_PFOCR_SETS:,}} PFOCR sets, {{NUM_MESH_TERMS:,}} MeSH terms), protein interactions ({{NUM_STRING_NEIGHBORS:,}} STRING neighbors), and single-cell expression ({{NUM_SINGLE_CELL_CLUSTERS:,}} cell clusters). Model selection identified a parsimonious combination of mouse phenotypes and MSigDB as the top performer (see below), producing {{NUM_GENETRAIT_REL_PP90:,}} gene-trait relationships with posterior probability above 90\% and {{NUM_GENETRAIT_REL_PP50:,}} above 50\%.
